\Needspace{6\baselineskip}
\begin{center}\hypertarget{native-258-article185}{}\pdfbookmark[2]{185.}{article-185}185.\end{center}
Der bisher behandelte Gegenstand könnte Anlass zu mehreren eleganten analytischen Untersuchungen geben. Darauf gehen wir hier nicht näher ein, damit wir uns nicht zu weit von unserem Vorhaben entfernen. Aus demselben Grund müssen wir die Darstellung der Verfahren, mit denen sich die numerische Rechnung auf einen leichter handhabbaren Algorithmus zurückführen lässt, einer anderen Gelegenheit vorbehalten. Eine einzige Bemerkung sei hier gestattet. Ist die Zahl der \SourceLeaf{259}vorgelegten Funktionen oder Gleichungen beträchtlich, so wird die Rechnung vor allem deshalb etwas mühsamer, weil die Koeffizienten, mit denen die ursprünglichen Gleichungen multipliziert werden müssen, um $P$, $Q$, $R$, $S$ usw. zu erhalten, meist unhandliche Dezimalbrüche enthalten. Erscheint es in diesem Fall nicht lohnend, diese Multiplikationen mithilfe von Logarithmentafeln so genau wie möglich auszuführen, genügt es in den meisten Fällen, anstelle dieser Multiplikatoren andere zu verwenden, mit denen es sich leichter rechnen lässt und die nur wenig von ihnen abweichen. Diese Vereinfachung kann nur dann merkliche Fehler verursachen, wenn das Maß der Präzision bei der Bestimmung der Unbekannten erheblich kleiner ausfällt als die Präzision der ursprünglichen Beobachtungen.
